% !TeX root = ../Cosmology.tex

% --------------------------------------------------------
\chapter{Квантові флуктуації і збурення}
\label{sec:quantum}
%% --------------------------------------------------------

У §\,\ref{sec:inflaton} інфлатон розглядався як класичне однорідне
поле $\bar{\phi}(t)$, а у §\,\ref{sec:metric} ми побудували математичний апарат
для опису збурень метрики. Тепер постає фундаментальне питання про
їхнє походження.

Будь-яке квантове поле неминуче флуктуює~--- навіть у вакуумному стані.
Ці нульові коливання можна представити як сукупність хвиль (мод Фур'є)
найрізноманітніших масштабів~--- від субмікроскопічних до космологічних,
де кожна мода характеризується своїм хвильовим числом $k$ та відповідною
фізичною довжиною хвилі $\lambda = 2\pi a(t)/k$. Під час інфляції ці
нульові коливання існують скрізь, зокрема й у нашій майбутній області простору.
Але стрімке інфляційне розширення розтягує фізичну довжину хвилі кожної моди
($\lambda \propto a(t)$). Щойно для моди з певним $k$ виконується умова
$\lambda \approx cH^{-1}$, вона перетинає горизонт Хаббла. Протилежні фази
цієї хвилі втрачають причинний контакт: вони більше не «взаємодіють» між
собою, і часові осциляції припиняються. Амплітуда хвиль класично
заморожується на тому значенні, яке мода мала в момент виходу з-під горизонту.

Проте горизонт Хаббла не залишається фіксованим назавжди. Після закінчення
інфляції та переходу до радіаційно-домінованої епохи закон розширення Всесвіту
кардинально змінюється. Якщо під час інфляції масштабний фактор зростав
прискорено ($\ddot{a} > 0$), то у гарячу епоху, коли домінує релятивістська
плазма, він сповільнюється як $a(t) \propto t^{1/2}$. Відповідно, параметр
Хаббла спадає як $H = \dot{a}/a = 1/(2t)$. Це призводить до того, що
співвідношення між сопутнім горизонтом Хаббла та масштабним фактором
змінюється: під час інфляції $d(aH)^{-1}/dt < 0$ (горизонт стискається
в сопутніх координатах), а після її закінчення
\begin{equation}
\frac{d}{dt}\left( \frac{1}{aH} \right) > 0.
\label{eq:horizon_growth}
\end{equation}
(Тут $c=1$ в природних одиницях, що використовуються в усій главі.)
На мові фізичних відстаней це означає, що реальний радіус горизонту
$R_H = 1/H \propto t$ зростає лінійно, випереджаючи ріст довжини хвилі
збурень $\lambda \propto t^{1/2}$. Внаслідок цієї нерівності раніше
заморожені моди одна за одною знову повертаються під горизонт~--- і саме тут
вони «розморожуються», знову починаючи осцилювати з тією самою амплітудою,
яку отримали під час інфляції. Ця амплітуда стає початковою умовою для
коливань гарячої фотон-баріонної плазми й зрештою відбивається в температурній
анізотропії реліктового випромінювання (CMB) та великомасштабній структурі Всесвіту.

Для того щоб кількісно описати цей складний процес і розрахувати конкретні
спостережувані наслідки, реальні збурення розкладають за типами симетрії
простору. Оскільки на ранніх етапах еволюції ці коливання залишаються вкрай
малими відхиленнями від однорідного фону (як ми побачимо далі, їхня амплітуда
становить усього близько однієї стотисячної долі від середньої густини,
$\delta\rho/\rho \sim 10^{-5}$), їхній розвиток можна розглядати ізольовано
в межах теорії збурень першого порядку.

У лінійному наближенні квантові збурення поділяються на два незалежні класи:
\begin{enumerate}
  \item \textbf{Флуктуації інфлатона} $\delta\phi(\eta, \mathbf{x}) = \phi(\eta, \mathbf{x}) - \bar{\phi}(\eta, \mathbf{x})$~---
        породжують первинні скалярні збурення кривини, що стають
        «насінням» для формування галактик та їхніх скупчень.
  \item \textbf{Флуктуації метрики} $h_{ij}(\eta, \mathbf{x})$~---
        тензорні збурення простору-часу, які інтерпретуються як первинні
        гравітаційні хвилі.
\end{enumerate}
Ці два класи еволюціонують незалежно і залишають кардинально різні відбитки
у поляризації реліктового випромінювання (так звані $E$- та $B$-моди, які ми
детально розглянемо у §\,\ref{sec:pol}).

%% --- Пропедевтична карта розділу ---
\bigskip
\noindent\textbf{Логічна карта розділу.}
Перш ніж заглибитись у математику, окреслимо маршрут~---
бо шлях від квантового вакууму до температурних плям CMB
неочевидний і має кілька несподіваних поворотів.

% ------------------------------
% Логічна карта розділу (збурення) - горизонтальна, детальна
% ------------------------------
\begin{center}
\begin{tikzpicture}[roadmap]

% Основні ноди
\node[box] (delta) {$\delta\phi_k$};
\node[box, right=of delta] (freeze) {Замерзання};
\node[box, right=of freeze] (R) {$\mathcal{R}_k$};
\node[box, right=of R] (Phi) {$\Phi_k$};
\node[box, right=of Phi, fill=red!5, draw=red!70!black] (dTT) {$\dfrac{\delta T}{T}$};

% Стрілки
\draw[arrow] (delta) -- (freeze);
\draw[arrow] (freeze) -- (R);
\draw[arrow] (R) -- (Phi);
\draw[arrow] (Phi) -- (dTT);

% Підписи зверху
\node[font=\tiny\bfseries, text=blue!70!black, above=0.1cm of delta] {Квантовий шум};
\node[font=\tiny\bfseries, text=blue!70!black, above=0.1cm of freeze] {Вихід за горизонт};
\node[font=\tiny\bfseries, text=blue!70!black, above=0.1cm of R] {Геометрична складка};
\node[font=\tiny\bfseries, text=blue!70!black, above=0.1cm of Phi, align=center] {Гравітаційний\\ потенціал};
\node[font=\tiny\bfseries, text=red!70!black, above=0.1cm of dTT] {Анізотропія CMB};

% Пояснення знизу
\node[font=\scriptsize, gray, below=0.1cm of delta] {осцилює під горизонтом};
\node[font=\scriptsize, gray, below=0.1cm of freeze] {$k = aH$};
\node[font=\scriptsize, gray, below=0.1cm of R, align=center] {калібрувально-інваріантний\\$\mathcal{R}_k = \mathrm{const}$};
\node[font=\scriptsize, gray, below=0.1cm of Phi] {$\Phi_k = -\frac{2}{3}\mathcal{R}_k$};
\node[font=\scriptsize, gray, below=0.1cm of dTT] {акустичні піки, плями};

\end{tikzpicture}
\end{center}

\medskip
\noindent\textbf{Крок 1: $\delta\phi_k$ як квантовий шум.}
Починаємо з найпростішого об'єкта~--- флуктуації самого поля
інфлатона $\delta\phi_k$. Це звичайний квантовий осцилятор у
просторі де-Сіттера: під горизонтом він коливається, за
горизонтом~--- замерзає з амплітудою $H/\sqrt{2k^3}$.
Саме тут народжується первинний шум.

\medskip
\noindent\textbf{Крок 2: чому не $\delta\rho_k$, а $\mathcal{R}_k$~--- і чому це зручно.}
Здавалося б, природна міра неоднорідності~--- це збурення
густини $\delta\rho_k$. Але $\delta\rho_k$ залежить від вибору
системи координат (калібрування): просто перейшовши в іншу
систему відліку, можна штучно «знищити» або «створити» збурення.
Натомість \emph{супутнє збурення кривини}
\begin{equation*}
  \mathcal{R}_k \equiv \Psi_k - \frac{H}{\dot{\bar\phi}}\,\delta\phi_k
\end{equation*}
є калібрувально-інваріантним: воно однакове в будь-якій системі
координат і описує реальну геометричну складку простору.
Головна перевага: за горизонтом $\mathcal{R}_k = \mathrm{const}$~---
воно «замерзає» і не змінюється ні під час розігріву, ні при
зміні рівняння стану. Це робить його ідеальним «конвертом»,
в якому первинна квантова інформація без втрат доставляється
від інфляційної епохи до моменту входу мод під горизонт у
гарячому Всесвіті.

\medskip
\noindent\textbf{Крок 3: від $\mathcal{R}_k$ до спостережень.}
Повернувшись під горизонт у \emph{радіаційно-домінованій} епосі,
$\mathcal{R}_k$ породжує ньютонівський гравітаційний потенціал
через співвідношення $\Phi_k = -\frac{2}{3}\mathcal{R}_k$
(коефіцієнт $2/3$ специфічний для радіаційної епохи; у
матеріально-домінованій він інший). Цей потенціал запускає
акустичні осциляції фотон-баріонної плазми~--- і саме їх ми
бачимо у вигляді піків на спектрі CMB і температурних плям на
карті Planck.

\medskip
\noindent Увесь цей ланцюжок~---
$\delta\phi_k \;\to\; \mathcal{R}_k \;\to\; \Phi_k \;\to\;
\delta T/T$~--- і є темою цього розділу.


%% --------------------------------------------------------
\section{Скалярні флуктуації у просторі де-Сіттера}
%% --------------------------------------------------------

Рівняння~\eqref{eq:klein_gordon} описує фонове поле $\bar{\phi}(t)$.
Лінеаризуємо його для малого збурення $\delta\phi = \phi - \bar{\phi}$,
нехтуючи членом $V''(\bar{\phi})\,\delta\phi$ у субхаббловому режимі
($k^2 \gg \mathcal{H}^2$, де масою поля можна знехтувати), і переходимо
до конформного часу $\eta$ та Фур'є-простору. Використовуючи
$\dot{f} = f'/a$ і $\ddot{f} = (f'' - \mathcal{H}f')/a^2$\footnote{%
З означення $d\eta = dt/a$ маємо $\dot{f} = f'/a$. Диференціюємо ще раз:
$\ddot{f} = \dfrac{d}{dt}\!\left(\dfrac{f'}{a}\right)
= \dfrac{f''}{a^2} - \dfrac{a'}{a^2}\dfrac{f'}{a}
= \dfrac{1}{a^2}(f'' - \mathcal{H}f')$.},
отримуємо для $k$-ї моди \cite{mukhanov_2005}:
\begin{equation}\label{eq:fourier_field}
  \delta\phi_k'' + 2\mathcal{H}\,\delta\phi_k' + k^2\,\delta\phi_k = 0,
\end{equation}
де штрих~--- похідна за конформним часом $\eta$, $k \equiv |\mathbf{k}|$~---
хвильове число у сопутній системі координат (фізична довжина хвилі
$\lambda_\mathrm{phys} = a/k$ росте з часом, тоді як сопутне $k$ залишається сталим).
Поведінка кожної моди визначається співвідношенням між $\lambda_\mathrm{phys}$
та горизонтом Хаббла $H^{-1}$:

\begin{itemize}
  \item \textbf{Субхаббловий режим ($k \gg \mathcal{H}$):}
        Для осцилюючої моди $\delta\phi_k' \sim k\,\delta\phi_k$, тому хабблівський член
        за порядком величини: $2\mathcal{H}\delta\phi_k' \sim \mathcal{H}k\,\delta\phi_k \ll k^2\delta\phi_k$,
        оскільки $\mathcal{H} \ll k$. Рівняння~\eqref{eq:fourier_field} зводиться до
        звичайного гармонічного осцилятора:
        \begin{equation*}
          \delta\phi_k'' + k^2\,\delta\phi_k \approx 0.
        \end{equation*}
        Це рівняння вільного поля у пласкому просторі Мінковського~--- таке саме,
        як для фотона або будь-якого іншого квантового поля. Його квантовий
        вакуумний стан добре відомий: це стан з мінімальною енергією, де кожна мода $k$
        знаходиться у нульових коливаннях з амплітудою $\sim 1/\sqrt{2k}$ (як амплітуда нульових коливань звичайного
        квантового осцилятора із частотою $\omega = k$ у природній системі одиниць $\hbar = c = 1$). У розширюваному просторі де-Сіттера цей стан
        називається \emph{вакуумом Банча-Девіса} \cite{bunch_davies_1978}~--- це
        природний вибір початкового стану: на достатньо малих масштабах ($k \gg \mathcal{H}$)
        кривина простору непомітна, і поле «не знає», що воно знаходиться в де-Сіттері.
        А отже, його вакуум збігається зі звичайним вакуумом Мінковського. Саме цей
        стан і є природною початковою умовою. З урахуванням масштабного фактора $a$
        нормування амплітуди набуває вигляду:
        \begin{equation*}
          |\delta\phi_k(\eta)| \sim \frac{1}{a(\eta) \sqrt{2k}}.
        \end{equation*}
        Множник $1/a$ має простий фізичний зміст: сопутне хвильове число $k$ не змінюється
        з часом, але фізична довжина хвилі $\lambda_\mathrm{phys} = a/k$ розтягується
        разом із простором. Енергія моди «розмазується» по більшому об'єму, і амплітуда
        спадає як $1/a$~--- аналогічно до того, як температура реліктового
        випромінювання спадає обернено пропорційно масштабному фактору.
  \item \textbf{Суперхаббловий режим ($k \ll \mathcal{H}$):}
        Градієнтний член $k^2\delta\phi_k$ стає нехтовно малим порівняно з
        $\mathcal{H}^2\delta\phi_k$, і рівняння~\eqref{eq:fourier_field} перетворюється
        з осцилятора на рівняння згасання (тертя):
        \begin{equation*}
          \delta\phi_k'' + 2\mathcal{H}\,\delta\phi_k' \approx 0.
        \end{equation*}
        Член $k^2$ — це «пружина» осцилятора, яка змушує моду коливатися. Коли мода
        виходить за горизонт, $k/\mathcal{H} \to 0$ і пружина «зникає». Фізично це означає,
        що протилежні фази хвилі більше не перебувають у причинному контакті через надто
        швидке розширення простору~--- хвиля замерзає. Розв'язок для домінуючої моди:
        $\delta\phi_k = \mathrm{const}$. Мода «виходить за горизонт» і заморожується \cite{hawking_1982}.

\begin{mathbox}{Аналітичний розв'язок суперхабблової моди}
    Щоб зрозуміти, чому з рівняння другого порядку виникає константа, роз\-в'я\-же\-мо суперхабблову систему
    $\delta\phi_k'' + 2\mathcal{H}\,\delta\phi_k' \approx 0$ аналітично. Зробимо заміну для швидкості зміни поля:
    $Y(\eta) \equiv \delta\phi_k'$, що дає рівняння першого порядку $Y' + 2\mathcal{H}Y = 0$, або:
    \begin{equation*}
      \frac{dY}{Y} = -2\mathcal{H}\,d\eta.
    \end{equation*}
    Під час інфляції конформний параметр Хаббла змінюється як $\mathcal{H} = -1/\eta$ (де конформний час
    $\eta < 0$ і прямує до нуля). Підставивши це, отримуємо $dY/Y = 2\,d\eta/\eta$. Інтегрування дає:
    \begin{equation*}
      \ln Y = 2\ln\eta + \text{const} \implies Y(\eta) = C_1 \eta^2.
    \end{equation*}
    Тепер, щоб знайти саме поле $\delta\phi_k$, проінтегруємо отриману швидкість ще раз:
    \begin{equation*}
      \delta\phi_k(\eta) = \int Y(\eta)\,d\eta = \int C_1 \eta^2\,d\eta = A\eta^3 + B,
    \end{equation*}
    де $A$ і $B$~--- сталі інтегрування. Перший доданок ($A\eta^3$) за умов $\eta \to 0$ стрімко згасає
    (\emph{загасаюча мода}), тому глибоко за горизонтом від усього розв'язку залишається лише чиста константа
    (\emph{домінуюча мода}): $\delta\phi_k \approx B = \mathrm{const}$.

    Конкретне значення цієї константи $B$ визначається в момент перетину горизонту ($k = \mathcal{H} = a_* H$),
    коли хвиля тільки-но замерзає. «Зшиваючи» цей результат із амплітудою швидких вакуумних коливань під
    горизонтом $|\delta\phi_k| \sim 1/(a\sqrt{2k})$ у точці перетину, де масштабний фактор дорівнює $a_* = k/H$,
    ми однозначно фіксуємо величину константи:
    \begin{equation*}
      B = \left. \frac{1}{a\sqrt{2k}} \right|_{a = a_*} = \frac{1}{\left(\dfrac{k}{H}\right)\sqrt{2k}} = \frac{H}{\sqrt{2k^3}}.
    \end{equation*}
    Таким чином, суперхабблова амплітуда поля є фіксованим «стоп-кадром» квантового шуму в мить його виходу за горизонт~\cite{guth_pi_1982}.
\end{mathbox}
\end{itemize}

%З урахуванням значення для замороженої
%амплітуди $|\delta\phi_k| \;\to\; \frac{H}{\sqrt{2k^3}}$, первинний безрозмірний спектр потужності флуктуацій поля  на суперхабблових масштабах записується як~\cite{mukhanov_chibisov_1981}:
%\begin{equation}\label{eq:field_power_spec}
%  \mathcal{P}_{\delta\phi}(k) \equiv \frac{k^3}{2\pi^2}|\delta\phi_k|^2
%  = \frac{k^3}{2\pi^2} \left( \frac{H}{\sqrt{2k^3}} \right)^{\!2}
%  = \left(\frac{H}{2\pi}\right)^{\!2}.
%\end{equation}

З урахуванням значення для замороженої
амплітуди $|\delta\phi_k| \;\to\; \frac{H}{\sqrt{2k^3}}$, ми можемо розрахувати головну статистичну характеристику
цього просторового квантового шуму~--- первинний \emph{спектр потужності}\footnote{Термін запозичений із теорії сигналів та спектрального аналізу. У математичній статистиці спектр потужності~--- це просто  розподіл середнього квадрата відхилення (дисперсії) випадкового поля за просторовими масштабами $k$.}
флуктуацій~\cite{mukhanov_chibisov_1981}.

Щоб побачити його математичний зміст, знайдемо повну дисперсію поля $\langle\delta\phi^2\rangle$,
підсумувавши внески від усіх можливих фур'є-мод у тривимірному просторі хвильових векторів:
\begin{equation*}
  \langle\delta\phi^2\rangle = \int_0^\infty \frac{d^3k}{(2\pi)^3} |\delta\phi_k|^2.
\end{equation*}
Перейдемо у сферичні координати ($d^3k = 4\pi k^2 dk$) та перепишемо інтеграл так,
щоб виділити інтегрування за одиничним логарифмічним інтервалом масштабів $d\ln k = dk/k$:
\begin{equation*}
  \langle\delta\phi^2\rangle = \int_0^\infty \frac{4\pi k^2 dk}{(2\pi)^3} |\delta\phi_k|^2
  = \int_0^\infty \left( \frac{k^3}{2\pi^2} |\delta\phi_k|^2 \right) \frac{dk}{k}
  = \int_0^\infty \mathcal{P}_{\delta\phi}(k) \, d\ln k.
\end{equation*}
Величина у дужках, за визначенням, і є безрозмірним спектром потужності $\mathcal{P}_{\delta\phi}(k)$.
Підставляючи сюди вираз для нашої замороженої на горизонті амплітуди, отримуємо фінальний результат:
\begin{equation}\label{eq:field_power_spec}
  \mathcal{P}_{\delta\phi}(k) \equiv \frac{k^3}{2\pi^2}|\delta\phi_k|^2
  = \frac{k^3}{2\pi^2} \left( \frac{H}{\sqrt{2k^3}} \right)^{\!2}
  = \left(\frac{H}{2\pi}\right)^{\!2}.
\end{equation}

Дивовижний результат формули~\eqref{eq:field_power_spec} полягає в тому, що вся
залежність від хвильового числа $k$ повністю скоротилася (геометричний множник $k^3$
ідеально компенсував об'ємний фактор $k^3$ у знаменнику амплітуди). Це означає, що
інфляція генерує \emph{масштабно-інваріантний} спектр флуктуацій~--- амплітуда космологічного шуму є практично однаковою на будь-яких просторових масштабах. Цей результат відомий як \emph{спектр Гар\-рі\-со\-на--Зельдовича} \cite{Harrison1970, Zeldovich1972}~--- перший теоретичний
передбачений вигляд первинного спектра збурень.

\begin{keybox}{Зауваження: де-Сіттерівська температура і термодинамічний зміст спектра}
Результат~\eqref{eq:field_power_spec} має значно глибший зміст, ніж просто вдале
алгебраїчне скорочення степенів. Простір де-Сіттера має горизонт подій, обмеження доступу
до інформації за яким, згідно з Гіббонсом і Гокінгом \cite{gibbons_hawking_1977},
призводить до виникнення теплового випромінювання вакуумних флуктуацій із характерною
температурою:
\begin{equation*}
  T_\mathrm{dS} = \frac{H}{2\pi}.
\end{equation*}
Таким чином, комбінація $H/(2\pi)$ у формулі~\eqref{eq:field_power_spec} є прямим
наслідком квантової термодинаміки горизонту. Оскільки флуктуації всіх масштабів
народжуються з однаковою ефективною «температурою» де-Сіттерівського вакууму,
ми й отримуємо спектр, який майже не залежить від $k$.
\end{keybox}

Проте тут слід зробити важливе фізичне уточнення щодо часу: під час реальної інфляції
параметр Хаббла не є строгою константою, а повільно зменшується у режимі повільного
скочування (\emph{slow-roll}). Через це значення температури $T_\mathrm{dS}$ (а отже,
й параметра $H$) у формулі спектра береться індивідуально для кожної моди~--- \textbf{у момент
її власного перетину горизонту ($k = aH$)}. Моди великих просторових масштабів (малі $k$)
виходять за горизонт раніше, коли величина $H$ була більшою, а моди малих масштабів
(великі $k$)~--- пізніше, коли інфляція вже наближалася до фіналу і $H$ трохи знизився.
Саме цей динамічний ефект порушує ідеальну масштабну інваріантність і призводить до того,
що первинний спектр реального Всесвіту набуває легкого «червоного нахилу» (спектральний
індекс $n_s \approx 0.965 \neq 1$), що блискуче підтверджується сучасними спостереженнями
реліктового випромінювання \cite{planck_2018}.

%% --------------------------------------------------------
\section{Збурення кривини і первинний спектр}
%% --------------------------------------------------------

Досі ми розглядали флуктуації скалярного поля $\delta\phi$ так, наче вони живуть
у фіксованому просторі де-Сіттера. Проте у реальному Всесвіті поле інфлатона є
головним джерелом гравітації: його густина енергії викривляє простір-час. Тому,
згідно з рівняннями Ейнштейна, будь-які квантові коливання поля $\delta\phi$ неминуче
«продавлюють» геометрію Всесвіту, миттєво генеруючи збурення метрики.

Оскільки поле і метрика жорстко зв'язані між собою гравітаційною взаємодією, розділити
їх у чистому вигляді неможливо. Щобільше, сама величина збурення поля $\delta\phi = \phi - \bar\phi$
залежить від того, як саме ми проведемо просторові перерізи (гіперповерхні) у чотиривимірному
просторі-часі. Зміна системи координат (калібрування) може штучно створювати або знищувати
збурення поля, що є суто математичним артефактом.

\subsection*{Фізичний механізм: як флуктуація поля народжує кривину}

Щоб зрозуміти, як квантовий шум поля трансформується в геометрію, уявімо інфлатон як «космічний годинник»,
що своїм значенням $\phi(t)$ відраховує час до кінця інфляції. Через квантові флуктуації
цей годинник у різних точках простору йде неузгоджено. Величина $\delta t = \delta\phi / \dot{\bar\phi}$
має фундаментальний сенс: це \emph{локальна затримка в часі} (або «запізнення» поля у конкретній точці
порівняно з середнім фоновим рухом).

Оскільки Всесвіт стрімко розширюється з темпом Хаббла $H$, у тих областях, де поле запізнилося ($\delta\phi > 0$),
інфляція триває трохи довше, і простір встигає локально розширитися в $e^{H \delta t}$ разів сильніше.
Це додаткове нерівномірне розтягнення простору геометрично означає лише одне~--- появу реальних
складок просторової кривини: $\mathcal{R} \sim H \delta t = \frac{H}{\dot{\bar\phi}}\delta\phi$.

%% --------------------------------------------------------
\subsection*{Калібрувально-інваріантний опис}
%% --------------------------------------------------------

Щоб позбутися координатної невизначеності та враховувати цей ефект на будь-яких масштабах,
нам потрібна \emph{калібрувально-інваріантна} величина, яка фізично об'єднує збурення поля та геометрії
в один монолітний об'єкт. У конформно-ньютонівському калібруванні таке \emph{супутнє збурення кривини\footnote{Не плутати зі скалярною кривиною Рімана $R$
з рівнянь Ейнштейна: $\mathcal{R}$ є безрозмірною амплітудою
збурення метрики, тоді як $R$ має розмірність
$[\text{довжина}^{-2}]$.}}
(comoving curvature perturbation) визначається як \cite{bardeen_1980}:
\begin{equation}\label{eq:R_definition}
  \mathcal{R}_k \equiv \Psi_k - \frac{H}{\dot{\bar\phi}}\,\delta\phi_k.
\end{equation}%

Рівняння руху для $\mathcal{R}_k$ у Фур'є-просторі має вигляд \cite{bardeen_1980} (рівняння Муханова-Сасакі):
\begin{equation}\label{eq:R_eom}
  \mathcal{R}_k'' + \frac{2z'}{z}\,\mathcal{R}_k' + k^2\mathcal{R}_k = 0, \qquad z \equiv \frac{a\dot{\bar\phi}}{H}.
\end{equation}

\begin{mathbox}{Рівняння Муханова-Сасакі}
    Рівняння Муханова-Сасакі є результатом лінеаризації системи рівнянь Ейнштейна
    $\delta G^\mu_\nu = 8\pi G \delta T^\mu_\nu$ першого порядку малості. Ми використовуємо загальну
    збурену метрику~\eqref{eq:perturbed_metric}, проте оскільки зараз розглядаються суто
    \emph{скалярні збурення} (генеровані полем інфлатона), тензорний член гравітаційних хвиль
    покладається рівним нулю ($h_{ij} = 0$). Крім того, за відсутності анізотропного тиску
    скалярні потенціали тотожно рівні між собою: $\Phi = \Psi$.

    У Фур'є-просторі геометрична ліва частина ($\delta G^\mu_\nu$) та польова права частина ($\delta T^\mu_\nu$)
    для такої метрики мають наступну структуру компонент:

    \begin{enumerate}
      \item \textbf{Часова компонента ($\delta G^0_0 = 8\pi G \delta T^0_0$) --- аналог рівняння Пуассона:}
      \begin{align*}
        \delta G^0_0 &= \frac{2}{a^2} \left[ k^2\Psi_k + 3\mathcal{H}(\Psi_k' + \mathcal{H}\Psi_k) \right], \\
        \delta T^0_0 &= -\delta\rho_k = - \frac{1}{a^2} \left[ \dot{\bar{\phi}}\delta\phi_k' - \dot{\bar{\phi}}^2\Psi_k - (\dot{\bar{\phi}}' + 2\mathcal{H}\dot{\bar{\phi}})\delta\phi_k \right].
      \end{align*}
      Тут ми врахували фонове рівняння руху поля~\eqref{eq:klein_gordon} ($a^2 V_{,\phi} = -\dot{\bar{\phi}}' - 2\mathcal{H}\dot{\bar{\phi}}$). Ця рівність показує, як квантові неоднорідності густини енергії поля ($\delta\rho_k$) миттєво продавлюють просторову геометрію ($k^2\Psi_k$).

      \item \textbf{Просторово-часова компонента ($\delta G^0_i = 8\pi G \delta T^0_i$) --- баланс імпульсу:}
      \begin{align*}
        \delta G^0_i &= \frac{2}{a^2} i k_i (\Psi_k' + \mathcal{H}\Psi_k), \\
        \delta T^0_i &= \frac{1}{a^2} i k_i \left( \dot{\bar{\phi}} \delta\phi_k \right).
      \end{align*}
      Після скорочення загального хвильового множника $ik_i$ це рівняння зводиться до прямого лінійного зв'язку між швидкістю зміни геометричного потенціалу та амплітудою поля: $\Psi_k' + \mathcal{H}\Psi_k = 4\pi G \dot{\bar{\phi}} \delta\phi_k$.
    \end{enumerate}

    Оскільки $\Psi_k$ та $\delta\phi_k$ диференціально переплетені у цих компонентах ЗТВ, для розплутування системи
    використовують визначення~\eqref{eq:R_definition}. З просторово-часової компоненти маємо:
    \begin{equation*}
      \Psi_k' + \mathcal{H}\Psi_k = 4\pi G\,\dot{\bar\phi}\,\delta\phi_k
      \;\Longrightarrow\;
      \Psi_k = \frac{\mathcal{H}}{\dot{\bar\phi}}\,\delta\phi_k + \mathcal{R}_k
    \end{equation*}
    (останній крок~--- з означення~\eqref{eq:R_definition}).
    Підставляючи це у часову компоненту і застосовуючи закон збереження
    $\nabla_\mu \delta T^{\mu\nu} = 0$, більшість членів скорочується
    завдяки фоновому рівнянню руху~\eqref{eq:klein_gordon},
    і система зводиться до одного рівняння для $\mathcal{R}_k$:
    \begin{equation*}
      \mathcal{R}_k'' + 2\left(\mathcal{H} + \frac{\ddot{\bar\phi}}{\dot{\bar\phi}} - \frac{\dot{H}}{H}\right)\mathcal{R}_k' + k^2\mathcal{R}_k = 0.
    \end{equation*}
    Громіздкий коефіцієнт «космічного тертя» у дужках математично тотожний логарифмічній похідній $2z'/z$ від функції
    $z \equiv \frac{a\dot{\bar\phi}}{H}$. Використовуючи раніше введені фонові характеристики інфлатона~\eqref{eq:inflaton_rho_p},
    цю функцію можна записати у красивому термодинамічному вигляді:
    \begin{equation*}
       z^2 = \frac{a^2 \dot{\bar{\phi}}^2}{H^2} = \frac{a^2 (\bar{\rho}_\phi + \bar{p}_\phi)}{H^2} = 3 M_{\mathrm{Pl}}^2 a^2 (1 + w_\phi),
    \end{equation*}
    де $w_\phi \equiv \bar{p}_\phi / \bar{\rho}_\phi$~--- параметр рівняння стану поля. Заміна дужок на $2z'/z$ і дає
    фінальне компактне рівняння Муханова~--- Сасакі~\eqref{eq:R_eom}.

    Тепер фізичний зміст $z$ стає очевидним: якщо Всесвіт заповнений чистою космологічною сталою ($w_\phi = -1$), то $z = 0$, і збурення кривини не генеруються, бо простір ідеально однорідний. Народження кривини $\mathcal{R}_k$ можливе лише тоді, коли $1 + w_\phi \neq 0$, тобто коли інфлатон має кінетичну енергію і повільно скочується, а сама функція $z$~--- це «гідродинамічно скоригований» масштабний фактор.
\end{mathbox}


%% --------------------------------------------------------
\subsection*{Еволюція розв'язків на різних масштабах: від кванту до класики}
%% --------------------------------------------------------

Рівняння~\eqref{eq:R_eom} спрощується канонічною заміною
$u_k \equiv z\mathcal{R}_k$, яка виносить «тертя» і зводить
задачу до осцилятора зі змінною масою:
\begin{equation}\label{eq:u_definition}
  u_k(\eta) \equiv z(\eta) \mathcal{R}_k(\eta).
\end{equation}
У Фур'є-просторі еволюція рівняння~\eqref{eq:R_eom} приймає вигляд:
\begin{equation}\label{eq:R_eom_u}
  u_k'' + \left( k^2 - \frac{z''}{z} \right) u_k = 0,
\end{equation}
де штрих означає похідну за конформним часом $\eta$. Поведінка розв'язків цього рівняння повністю визначається просторовим масштабом хвилі відносно причинного горизонту Хаббла $\mathcal{H} = aH$. Залежно від співвідношення між геометричним та кінетичним членами у дужках, виділяють два фундаментальних режими еволюції збурень:

\begin{enumerate}
\item \textbf{Глибоко субхаббловий режим ($k \gg \mathcal{H}$ або $\lambda \ll H^{-1}$):}
  На ранніх етапах інфляції довжина хвилі збурення $\lambda$ є мізерною порівняно з радіусом причинності. У цьому режимі конформний фактор гравітаційного потенціалу у рівнянні Муханова--Сасакі є нехтовно малим порівняно з кінетичним членом: $k^2 \gg |z''/z| \approx 2/\eta^2$. За цієї умови загальне рівняння руху для квантової змінної $u_k$ спрощується до вигляду:
  \begin{equation}\label{eq:MS_sub_horizon}
    u_k'' + k^2 u_k \approx 0.
  \end{equation}
  Рівняння~\eqref{eq:MS_sub_horizon} за своєю структурою тотожне рівнянню звичайного гармонічного осцилятора у пласкому просторі-часі Мінковського. Його фундаментальним квантовим розв'язком, що задовольняє умове мінімуму енергії, є вакуумний стан Банча--Дейвіса:
  \begin{equation}\label{eq:bunch_davies}
    u_k(\eta) = \frac{1}{\sqrt{2k}} e^{-ik\eta}.
  \end{equation}
  У цьому режимі комплексна фаза хвилі стрімко обертається, а її квантове середнє значення дорівнює нулю.

  Метричні потенціали $\Psi_k$ та $\Phi_k$ тут є нехтовно малими. З лінеаризованих рівнянь Айнштайна (аналога рівняння Пуассона) випливає, що високочастотні коливання інфлатона генерують потенціал амплітудою $\Psi_k \sim \dot{\bar{\phi}}\delta\phi_k / (k M_{\mathrm{Pl}}^2)$. Пряме порівняння метричного та польового внесків у повному релятивістському визначенні кривизни показує їхнє відношення:
  \begin{equation}\label{eq:psi_suppression}
    \frac{\Psi_k}{\frac{\mathcal{H}}{\dot{\bar{\phi}}}\delta\phi_k} \sim \frac{\dot{\bar{\phi}}^2}{M_{\mathrm{Pl}}^2 H^2} \cdot \frac{\mathcal{H}}{k} \sim 2\varepsilon \left(\frac{\mathcal{H}}{k}\right) \ll 1.
  \end{equation}
  Формула~\eqref{eq:psi_suppression} демонструє подійне придушення гравітаційного відгуку: по-перше, через малість параметра повільного скочування ($\varepsilon \ll 1$), а по-друге~--- геометрично, оскільки мікроскопічна хвиля ($k \gg \mathcal{H}$) не встигає суттєво викривити простір на своїх масштабах. Фізично простір-час здається такій моді практично пласким. Тому, зважаючи на кінематичний зв'язок $u_k \approx a\,\delta\phi_k$, збурення супутньої кривизни повністю зводиться до чистої коливальної вакуумної динаміки самого поля інфлатона \cite{mukhanov_2005}:
  \begin{equation}\label{eq:curvature_perturbation_sub}
    \mathcal{R}_k \approx -\frac{\mathcal{H}}{\dot{\bar{\phi}}}\,\delta\phi_k = -\frac{H}{\dot{\bar{\phi}}}\,\delta\phi_k.
  \end{equation}

\item \textbf{Суперхаббловий режим ($k \ll \mathcal{H}$ або $\lambda \gg H^{-1}$):}
  Оскільки Всесвіт розширюється експоненціально, фізична довжина хвилі збурення зростає значно швидше за майже сталий радіус Хаббла $c/H$. Мода перетинає горизонт ($k \approx \mathcal{H}$) і виходить за його межі. У суперхаббловому режимі член просторового градієнта $k^2\mathcal{R}_k \to 0$. Якщо знехтувати цим членом, рівняння Муханова--Сасакі легко інтегрується в загальному вигляді:
  \begin{equation}\label{eq:R_integration}
    \mathcal{R}_k(\eta) = C_1 + C_2 \int \frac{d\eta}{z^2(\eta)},
  \end{equation}
  де супутній параметр $z \equiv a \dot{\bar{\phi}} / \mathcal{H}$. У наближенні повільного скочування швидкість поля $\dot{\bar{\phi}}$ та параметр Хаббла $H$ змінюються дуже повільно, тому $z(\eta) \propto a(\eta)$. Згадавши, що для де-Сіттерівського розширення масштабний фактор пов'язаний із конформним часом як $a(\eta) = -1/(H\eta)$, де $\eta \in (-\infty, 0)$, ми можемо явно обчислити цей інтеграл:
  \begin{equation}\label{eq:R_integration_explicit}
    \int \frac{d\eta}{z^2} \propto \int \frac{d\eta}{a^2} = \int H^2 \eta^2 d\eta = \frac{H^2 \eta^3}{3} \propto \frac{1}{a^3}.
  \end{equation}
  Підставляючи результат~\eqref{eq:R_integration_explicit} у загальний розв'язок~\eqref{eq:R_integration}, остаточно знаходимо часову еволюцію моди поза горизонтом:
  \begin{equation}\label{eq:R_super_horizon_sol}
    \mathcal{R}_k(\eta) = C_1 (k) + C_2 (k) \cdot \frac{1}{a^3(\eta)}.
  \end{equation}

    Отже, поза горизонтом залишається лише стала мода~---
    $\mathcal{R}_k$ буквально «фотографує» квантовий шум
    у момент виходу за горизонт і зберігає цей знімок назавжди.

  Математичний вигляд розв'язку~\eqref{eq:R_super_horizon_sol} має фундаментальний фізичний зміст, який кодує в собі поведінку збурень:
  \begin{itemize}
    \item \textbf{Загасаюча мода ($\propto a^{-3}$):} Другий член, що походить від конформного інтеграла, стрімко руйнується експоненціальним розширенням простору. Всього за кілька $e$-folds інфляції ця мода повністю обнуляється, не залишаючи жодного сліду в еволюції Всесвіту.
    \item \textbf{Зростаюча (стаціонарна) мода ($C_1 = \mathrm{const}$):} Перший член абсолютно не залежить від часу. Це означає, що амплітуда хвилі буквально \emph{«заморожується»} на тому рівні, який вона мала в момент перетину горизонту Хаббла ($k = aH$).
  \end{itemize}

    Фізично це замерзання означає, що вершина і низина
    гігантської хвилі рознесені у просторі на відстань,
    більшу за радіус причинності $c/H$. Локальні області
    простору більше не можуть обмінюватися світловими чи
    звуковими сигналами, а гідродинамічний тиск плазми
    принципово не здатний перерозподілити енергію та згладити
    збурення. Коливання повністю припиняються: комплексна
    фаза хвилі застигає. Квантова мода опиняється у сильно
    стиснутому стані, в якому некомутаційні поправки
    $[\hat{\mathcal{R}}_k, \hat{\mathcal{R}}_k'] \sim \hbar$
    стають нехтовно малими порівняно з класичною амплітудою~---
    і мода ефективно поводиться як звичайне
    \emph{класичне стохастичне гауссове поле}.

  Стабільність $\mathcal{R}_k = \mathrm{const}$ поза горизонтом робить її абсолютно невразливою до будь-яких мікрофізичних процесів наприкінці інфляції, зо\-крема до радикальної зміни рівняння стану речовини під час розігріву (reheating) Всесвіту. \emph{Це перетворює $\mathcal{R}_k$ на ідеальну «сейфову» початкову умову, в якій Всесвіт без жодних втрат чи зашумлення переносить первинну інформацію від квантової епохи ($t \sim 10^{-36}$ с) до моменту входу в гарячу фазу.}
\end{enumerate}

\begin{keybox}{Геометричний зв'язок Простору і Часу}
Важливо уникати термінологічного ступору щодо природи збурення супутньої кривизни $\mathcal{R}_k$. Безрозмірна величина $\mathcal{R}_k$ сидить у \emph{просторовій} частині метрики ($g_{ij} \approx a^2(1-2\mathcal{R})\delta_{ij}$), безпосередньо описуючи локальне розтягнення чи стиснення координатної сітки, що рухається разом із плазмою.

Математичний зв'язок цієї величини зі справжньою геометричною кривизною розкривається через розрахунок тривимірного скаляра кривизни $^{(3)}R$ у першому порядку теорії збурень:
\begin{equation}\label{eq:spatial_curvature}
  ^{(3)}R = \frac{4}{a^2} \nabla^2 \mathcal{R}
  \quad\xrightarrow{\text{Фур'є}}\quad
  ^{(3)}R_k = -\frac{4k^2}{a^2} \mathcal{R}_k.
\end{equation}
Рівняння~\eqref{eq:spatial_curvature} показує, що $\mathcal{R}_k$ є просторовим «гравітаційним потенціалом», лапласіан якого ($\nabla^2 \equiv -k^2$) задає реальне викривлення тривимірного простору.

Проте в ЗТВ простір і час жорстко зчеплені рівняннями Айнштайна ($G_{\mu\nu} = 8\pi G T_{\mu\nu}$). Тому просторова деформація геометрії поза горизонтом дзеркально відбивається у \emph{часову} компоненту метрики ($g_{00} \approx -(1+2\Phi)$), миттєво генеруючи класичний ньютонівський гравітаційний потенціал $\Phi_k$. На початковому етапі радіаційно-домінованої епохи ($w=1/3$) цей релятивістський зв'язок є строго фіксованим:
\begin{equation}\label{eq:R_to_Phi}
  \Phi_k = -\frac{2}{3}\mathcal{R}_k.
\end{equation}
Знак «мінус» має глибокий фізичний сенс: там, де простір під час інфляції був сильніше розтягнутий ($\mathcal{R}_k > 0$), виникає класична гравітаційна яма ($\Phi_k < 0$). Щойно після закінчення інфляції відповідна мода знову увійде під горизонт, ця часова яма почне притягувати речовину, змушуючи фотон-баріонне желе стискатися і запускаючи акустичні осциляції.
\end{keybox}

Оскільки $\mathcal{R}_k$ є класичною випадковою гауссовою величиною, ми описуємо її через статистичні характеристики. За визначенням, безрозмірний спектр потужності випадкового поля пов'язаний із середнім квадратом його Фур'є-мод як $\mathcal{P}_X(k) \equiv \frac{k^3}{2\pi^2} \langle |X_k|^2 \rangle$. Домножуючи співвідношення~\eqref{eq:curvature_perturbation_sub} на коефіцієнт $\frac{k^3}{2\pi^2}$ в момент перетину горизонту ($k = aH$), отримуємо прямий зв'язок між первинними спектрами поля та кривизни:
\begin{equation}\label{eq:spec_relation}
  \mathcal{P}_{\mathcal{R}}(k) = \left(\frac{H}{\dot{\bar{\phi}}}\right)^{\!2} \mathcal{P}_{\delta\phi}(k).
\end{equation}
Підставляючи сюди спектр квантових флуктуацій вільного скалярного поля у просторі де Сіттера $\mathcal{P}_{\delta\phi}(k) = (H/2\pi)^2$, остаточно знаходимо фундаментальну амплітуду первинних збурень кривизни:
\begin{equation}\label{eq:curvature_power_spec}
  \mathcal{P}_{\mathcal{R}}(k) = \left. \frac{H^4}{4\pi^2\dot{\bar{\phi}}^2} \right|_{k=aH}.
\end{equation}

Відхилення від суворо плоського (масштабно-інваріантного) спектру Гар\-рі\-со\-на-Зельдовича параметризується спектральним індексом $n_s$:
\begin{equation}\label{eq:ns}
  n_s - 1 \equiv \frac{d\ln\mathcal{P}_{\mathcal{R}}}{d\ln k} = -2\varepsilon - \eta_V,
\end{equation}
де $\varepsilon$ та $\eta_V$~--- параметри повільного скочування \eqref{eq:eps} та \eqref{eq:eps_eta_V}, обчислені в момент виходу моди $k$ за горизонт.

\begin{keybox}{Математичне виведення співвідношення для $n_s$}
Оскільки мода перетинає горизонт при $k = aH \approx \mathrm{const} \cdot a$, то диференціал хвильового числа прямо пов'язаний із числом $e$-folds: $d\ln k \approx d\ln a = dN$. Тоді диференціюванням рівняння~\eqref{eq:curvature_power_spec} за числом $e$-folds маємо:
\begin{equation*}
  n_s - 1 = \frac{d}{dN}\ln\! \left(\frac{H^4}{\dot{\bar\phi}}^2\right) = 4\frac{\dot H}{H^2} - 2\frac{\ddot{\bar\phi}}{H\dot{\bar\phi}} = -4\varepsilon + 2\eta_\mathrm{H}.
\end{equation*}
Диференціюванням за часом рівняння класичного атрактора повільного скочування $3H\dot{\bar\phi} \approx -V'(\bar\phi)$ та враховуючи поправки другого порядку \cite{mukhanov_2005}, знаходимо зв'язок хабблівського параметра другого порядку з параметрами потенціалу: $\eta_\mathrm{H} \approx \varepsilon - \tfrac{1}{2}\eta_V$. Підставляючи це значення, остаточно отримуємо:
\begin{equation*}
  n_s - 1 = -4\varepsilon + 2\bigl(\varepsilon - \tfrac{1}{2}\eta_V\bigr) = -2\varepsilon - \eta_V. \qquad\square
\end{equation*}
\end{keybox}

Космічна обсерваторія \textit{Planck} 2018 виміряла це значення з безпрецедентною точністю: $n_s = 0.9649 \pm 0.0042$, що строго підтверджує динаміку повільного кочення (спектр є «червоним», тобто нахиленим, $n_s < 1$, що виключає статичні моделі) \cite{planck_2018}.

Проте, якщо на шляху інфлатона трапляються локальні аномалії потенціалу (наприклад, різкі зміни нахилу чи сходинки), швидкість кочення $\dot{\bar\phi}$ може на короткий час стрімко впасти. Оскільки первинний спектр~\eqref{eq:curvature_power_spec} обернено пропорційний квадрату швидкості ($\mathcal{P}_{\mathcal{R}} \propto \dot{\bar{\phi}}^{-2}$), таке гальмування викликає локальний сплеск амплітуди флуктуацій, формуючи так званий \textbf{зламаний спектр}. Фізичні наслідки та феноменологічні обмеження таких аномалій ми детально розберемо у §\,\ref{sec:future_research}.

Повернувшись назад під причинний горизонт Хаббла вже у постінфляційну гарячу епоху, ці стабільні класичні моди $\mathcal{R}_k$ миттєво «вмикають» градієнти тиску і стають тими самими когерентними початковими умовами у фазі косинуса для гідродинаміки фотон-баріонної плазми, які ми детально вивчили у розділі про акустичні осциляції (BAO).

%% --------------------------------------------------------
\section{Динаміка первинних гравітаційних хвиль}
%% --------------------------------------------------------

З'ясуємо, як еволюціонують тензорні збурення $h_{ij}$. Підставляючи тензорну частину збуреної метрики у лінеаризовані рівняння Ейнштейна, отримуємо хвильове рівняння для кожної Фур'є-моди:
\begin{equation}\label{eq:tensor_wave_eq}
  h_{ij}'' + 2\mathcal{H}\,h_{ij}' + k^2\,h_{ij} = 0.
\end{equation}
Його структура ідентична структурі рівнянь для скалярного поля (осцилятор із хабблівським тертям). Тому тензорні моди так само заморожуються за горизонтом, а їхній первинний спектр визначається виключно масштабом Хаббла під час інфляції \cite{starobinsky_1979}:
\begin{equation}\label{eq:tensor_power_spec}
  \mathcal{P}_h(k) = \frac{2}{\pi^2 M_\mathrm{Pl}^2} \left(\frac{H}{2\pi}\right)^{\!2} = \frac{H^2}{\pi^2 M_\mathrm{Pl}^2}.
\end{equation}
Присутність планківської маси $M_\mathrm{Pl}^{-2}$ підкреслює квантово-гравітаційну природу цих хвиль та слабкість їхньої взаємодії. Експериментально амплітуда тензорної моди шукається через $B$-моди поляризації реліктового випромінювання за допомогою тензорно-скалярного відношення $r$:
\begin{equation}\label{eq:tensor_scalar_ratio}
  r \equiv \frac{\mathcal{P}_h(k_*)}{\mathcal{P}_{\mathcal{R}}(k_*)} = 16\varepsilon.
\end{equation}
Сучасне обмеження $r < 0.036$ \cite{bicep_keck_2021} жорстко лімітує енергетичний масштаб верхніх меж інфляції, відсікаючи великий клас теоретичних моделей. Детальний аналіз зв'язку між $r$ та поляризацією наведено у §\,\ref{sec:pol}.

\begin{keybox}{Зауваження про квантово-класичний перехід скалярних і тензорних мод}
Після виходу за горизонт обидва типи мод~--- скалярні
$\mathcal{R}_k$ і тензорні $h_{ij}$~--- поводяться як
класичні стохастичні поля. Практичним механізмом є
\emph{стиснення квантового стану}: у просторі де-Сіттера
дисперсія амплітудної квадратури зростає як $e^{2Ht}$,
тоді як спряжена фазова компонента пригнічується як
$e^{-2Ht}$. У такому граничному стані мода практично
невідрізняється від класичного стохастичного поля~---
незалежно від того, чи відбулася повноцінна декогеренція
\cite{polarski_starobinsky_1996, kiefer_polarski_1998}.

Точний механізм переходу (вибір базису вказівників, роль
довкілля) залишається відкритим концептуальним питанням.
Проте це питання концептуальне, а не вимірювальне: гауссова
статистика заморожених мод прекрасно узгоджується зі
спостереженнями~--- зокрема з відсутністю первинної
негауссовості за даними Planck~2018~\cite{planck_2018}.
\end{keybox}

\clearpage
% ============================================================
% Задачі до розділу 4: Квантові флуктуації і збурення
\section*{Задачі}
\addcontentsline{toc}{section}{Задачі}
% ============================================================

\begin{problem}{Вихід моди за горизонт та заморожування амплітуди}
Розглянемо квантову моду $\delta\phi_k$ скалярного поля в де-Сіттерівському просторі.

\begin{enumerate}
    \item Використовуючи рівняння \eqref{eq:fourier_field} та наближення
          $\mathcal{H} = -1/\eta$ (де $\eta < 0$), покажіть, що в суперхаббловому
          режимі ($k \ll \mathcal{H}$) розв'язок має вигляд $\delta\phi_k = A + B\eta^3$.
    \item Чому домінуючою є мода $\delta\phi_k = \mathrm{const}$? Поясніть фізично,
          чому амплітуда заморожується після перетину горизонту.
    \item Оцініть момент виходу моди за горизонт $k = aH$ через конформний час $\eta$.
          Який зв'язок між $k$ та $\eta_*$ у точці перетину?
\end{enumerate}
\end{problem}

% -------------------------------------------------------------------

\begin{problem}{Спектр потужності та де-Сіттерівська температура}
Спектр потужності флуктуацій інфлатона на суперхабблових масштабах дорівнює
$\mathcal{P}_{\delta\phi}(k) = (H/2\pi)^2$.

\begin{enumerate}
    \item Поясніть, чому цей спектр називається \emph{масштабно-інваріантним}.
          Якби залежність була $\mathcal{P}_{\delta\phi} \propto k^{n_s-1}$, то яке
          значення $n_s$ відповідає масштабній інваріантності?
    \item Де-Сіттерівська температура Гіббонса--Гокінга дорівнює $T_{\mathrm{dS}} = H/(2\pi)$.
          Як цей факт пов'язаний із спектром $\mathcal{P}_{\delta\phi}$?
    \item Чому в реальному Всесвіті $n_s \approx 0.965 \neq 1$? Яка фізична причина
          порушення масштабної інваріантності?
\end{enumerate}
\end{problem}

% -------------------------------------------------------------------

\begin{problem}{Збурення кривини та його збереження за горизонтом}
Супутнє збурення кривини визначається як $\mathcal{R}_k = \Psi_k - \frac{H}{\dot{\bar\phi}}\delta\phi_k$.

\begin{enumerate}
    \item На субхабблових масштабах ($k \gg \mathcal{H}$) можна знехтувати $\Psi_k$.
          Чому це наближення виправдане? Як спрощується $\mathcal{R}_k$ у цьому випадку?
    \item У суперхаббловому режимі ($k \ll \mathcal{H}$) доведіть, що $\mathcal{R}_k = \mathrm{const}$,
          використовуючи рівняння Муханова-Сасакі \eqref{eq:R_eom} та той факт, що $z \propto a$ зростає.
    \item Чому збереження $\mathcal{R}_k$ за горизонтом є критично важливим для
          зв'язку інфляції зі спостереженнями CMB?
\end{enumerate}
\end{problem}

% -------------------------------------------------------------------

\begin{problem}{Тензорні збурення та параметр $r$}
Тензорні збурення (первинні гравітаційні хвилі) мають спектр потужності
$\mathcal{P}_h = H^2/(\pi^2 M_{\mathrm{Pl}}^2)$, а тензорно-скалярне відношення
визначається як $r = \mathcal{P}_h / \mathcal{P}_{\mathcal{R}}$.

\begin{enumerate}
    \item Використовуючи $\mathcal{P}_{\mathcal{R}} = H^4/(4\pi^2\dot{\bar\phi}^2)$ та
          $\varepsilon = \dot{\phi}^2/(2H^2 M_{\mathrm{Pl}}^2)$, доведіть, що $r = 16\varepsilon$.
    \item Сучасне обмеження $r < 0.036$ (BICEP/Keck + Planck). Яку верхню межу
          на енергетичний масштаб інфляції $V^{1/4}$ це дає? (Підказка: $H^2 \approx V/(3M_{\mathrm{Pl}}^2)$.)
    \item Чому виявлення $B$-мод поляризації CMB, викликаних $r \gtrsim 0.001$,
          вважається «святим граалем» космології?
\end{enumerate}
\end{problem}

% -------------------------------------------------------------------

\begin{problem}{Квантово-класичний перехід та стиснені стани}

У просторі де-Сіттера кожна суперхабблівська мода еволюціонує у \emph{стиснутий
когерентний стан}. Для такої моди дисперсія однієї квадратурної компоненти
зростає як $e^{2Ht}$, а спряженої — спадає як $e^{-2Ht}$.

\begin{enumerate}
    \item Чому зростання дисперсії однієї з квадратур дозволяє говорити про
          \emph{класичну} поведінку моди? (Підказка: згадайте співвідношення
          невизначеностей Гайзенберга.)
    \item Чи є цей перехід \emph{декогеренцією} у звичайному сенсі? Чому автори
          \cite{polarski_starobinsky_1996} стверджують, що стиснення саме по собі
          вже достатнє для класичності?
    \item Який експериментальний факт підтверджує, що збурення справді поводяться
          як класичне стохастичне поле, а не як квантова суперпозиція?
\end{enumerate}
\end{problem}